\chapter*{Ringraziamenti}
\addcontentsline{toc}{chapter}{Ringraziamenti}

Devo dire, prima di ogni altra cosa, di essere sempre stato più capace di esprimere gratitudine con i gesti e con la presenza che con le parole. Che queste poche righe valgano allora per qualcosa che spero di mostrare più pienamente nel modo in cui sto con le persone che vi compaiono.

Il mio primo ringraziamento va al mio relatore, Jonathan, che ha accolto un'idea eclettica e non convenzionale e l'ha lasciata crescere. Significa molto per me che sia un astrofisico che insegna machine learning: la prova vivente che si può essere profondamente competenti senza essere ristretti, e che alla curiosità è concesso muoversi anche di lato, non solo in profondità.

Ai miei genitori, il cui sostegno è stato infinito da quando ho memoria: mi avete dato la curiosità, il metodo e la libertà di diventare chi sono. Quasi tutto ciò a cui tengo, l'ho visto prima in voi.

A mio fratello, che non vive più con me ma a cui mi sento vicino lo stesso. Sei sempre stato per me una fonte di ispirazione --- per la tua empatia, la tua bontà non comune e, soprattutto, la tua disciplina. Non riesco a immaginare la mia vita senza la tua presenza.

Ai miei nonni, che hanno contribuito a crescermi e hanno trasmesso amore e compassione come se fossero le cose più naturali da insegnare. Sono immensamente grato di avervi avuto nella mia vita, e spero di portare avanti il vostro esempio.

Ai miei amici, che mi sono stati accanto in questi anni e hanno abbracciato un percorso e una visione che vanno controcorrente. Camminare su una strada non convenzionale è molto più leggero con persone che scelgono di percorrerne un tratto accanto a te.

Alle persone che mi ispirano ogni giorno senza saperlo, e che non ho mai incontrato: mi ricordate che cosa significhi pensare in grande, continuare a imparare e provare a contare qualcosa.

Ai miei compagni di corso, a tutti coloro che hanno preso parte allo studio e, più in generale, alle persone che hanno arricchito la mia vita in questi ultimi anni: questo lavoro porta in sé qualcosa di tutti voi.

Infine, ma non da ultimo, voglio ringraziare me stesso: per la tenacia di resistere nei momenti difficili e nei dubbi su di me, per aver tenuto fede ai miei valori, per le volte in cui sono riuscito a ispirare gli altri e per aver continuato a spingermi verso il futuro che voglio costruire, uno in cui possa diventare pienamente me stesso e contribuire a rendere il mondo un posto migliore, per gli esseri umani e non umani allo stesso modo.

\cleardoublepage
